\chapter*{Glossário}
\addcontentsline{toc}{chapter}{Glossário}
%================================================

\section*{Como utilizar este glossário}

Ao longo desta apostila foram apresentados conceitos provenientes de diversas
áreas do conhecimento, incluindo ecologia, biogeografia, climatologia,
estatística, ciência de dados espaciais e modelagem computacional.

Este glossário reúne os principais termos, siglas e conceitos utilizados em
Modelagem de Distribuição de Espécies (SDM), funcionando como uma referência
rápida para consulta durante a leitura, desenvolvimento de projetos e
interpretação dos resultados.

\begin{longtable}{
		p{0.18\textwidth}
		p{0.72\textwidth}
	}
	
	\caption{Glossário dos principais termos utilizados em Modelagem de Distribuição de Espécies.}
	\label{tab:glossario_sdm}
	\\
	
	\toprule
	\textbf{Termo} &
	\textbf{Definição}
	\\
	\midrule
	\endfirsthead
	
	\toprule
	\textbf{Termo} &
	\textbf{Definição}
	\\
	\midrule
	\endhead
	
	SDM &
	Species Distribution Model. Modelo utilizado para estimar a distribuição espacial de uma espécie em função de variáveis ambientais.
	\\
	
	ENM &
	Ecological Niche Model. Abordagem voltada à modelagem do nicho ecológico da espécie.
	\\
	
	HSM &
	Habitat Suitability Model. Modelo de adequabilidade de habitat.
	\\
	
	AUC &
	Area Under the Curve. Métrica baseada na curva ROC utilizada para avaliar a capacidade discriminatória do modelo.
	\\
	
	ROC &
	Receiver Operating Characteristic. Curva utilizada para avaliar o desempenho classificatório dos modelos.
	\\
	
	TSS &
	True Skill Statistic. Métrica que combina sensibilidade e especificidade.
	\\
	
	Kappa &
	Coeficiente de concordância utilizado para avaliar desempenho preditivo.
	\\
	
	Boyce Index &
	Métrica que avalia a concordância entre adequabilidade prevista e registros observados.
	\\
	
	VIF &
	Variance Inflation Factor. Indicador utilizado para detectar multicolinearidade entre variáveis preditoras.
	\\
	
	PCA &
	Principal Component Analysis. Técnica de redução de dimensionalidade baseada em componentes principais.
	\\
	
	GLM &
	Generalized Linear Model. Modelo linear generalizado amplamente utilizado em SDM.
	\\
	
	GAM &
	Generalized Additive Model. Extensão não linear dos modelos lineares generalizados.
	\\
	
	RF &
	Random Forest. Algoritmo baseado em múltiplas árvores de decisão.
	\\
	
	BRT &
	Boosted Regression Trees. Método baseado em árvores de regressão combinadas por boosting.
	\\
	
	Maxent &
	Maximum Entropy Modeling. Algoritmo amplamente utilizado para dados de presença apenas.
	\\
	
	Ensemble &
	Combinação de múltiplos modelos para produzir uma predição consensual.
	\\
	
	Background &
	Conjunto de condições ambientais utilizadas para representar a disponibilidade ambiental na área de calibração.
	\\
	
	Pseudo-ausência &
	Local sem registro da espécie utilizado para representar ausências em determinados algoritmos.
	\\
	
	Área M &
	Região acessível à espécie ao longo de sua história biogeográfica segundo o framework BAM.
	\\
	
	BAM &
	Framework conceitual que integra componentes Bióticos, Abióticos e de Movimento (dispersão).
	\\
	
	Nicho Fundamental &
	Conjunto de condições ambientais sob as quais a espécie pode sobreviver na ausência de restrições bióticas.
	\\
	
	Nicho Realizado &
	Porção do nicho fundamental efetivamente ocupada pela espécie.
	\\
	
	Adequabilidade Ambiental &
	Grau de compatibilidade entre as condições ambientais de um local e o nicho estimado da espécie.
	\\
	
	Transferibilidade &
	Capacidade do modelo manter desempenho quando aplicado a novas regiões ou períodos temporais.
	\\
	
	Extrapolação &
	Predição realizada sob condições ambientais não observadas durante a calibração.
	\\
	
	MESS &
	Multivariate Environmental Similarity Surface. Ferramenta para detectar extrapolação ambiental.
	\\
	
	MOP &
	Mobility-Oriented Parity. Método para identificar novidade ambiental em projeções espaciais ou temporais.
	\\
	
	CMIP6 &
	Coupled Model Intercomparison Project Phase 6. Conjunto internacional de modelos climáticos globais.
	\\
	
	GCM &
	Global Climate Model. Modelo climático global utilizado em projeções futuras.
	\\
	
	SSP &
	Shared Socioeconomic Pathway. Cenários socioeconômicos utilizados nas projeções climáticas.
	\\
	
	SSP126 &
	Cenário de baixas emissões e forte mitigação climática.
	\\
	
	SSP585 &
	Cenário de altas emissões e forte aquecimento global.
	\\
	
	LIG &
	Last Interglacial. Último Interglacial, aproximadamente 130 mil anos antes do presente.
	\\
	
	LGM &
	Last Glacial Maximum. Último Máximo Glacial, aproximadamente 21 mil anos antes do presente.
	\\
	
	Mid-Holocene &
	Holoceno Médio, aproximadamente 6 mil anos antes do presente.
	\\
	
	WorldClim &
	Base global de variáveis climáticas amplamente utilizada em SDM.
	\\
	
	CHELSA &
	Conjunto de variáveis climáticas de alta resolução derivadas de modelos atmosféricos.
	\\
	
	FAIR &
	Findable, Accessible, Interoperable and Reusable. Princípios internacionais para gestão de dados científicos.
	\\
	
	DOI &
	Digital Object Identifier. Identificador digital permanente para publicações e conjuntos de dados.
	\\
	
	Git &
	Sistema de controle de versão amplamente utilizado em ciência computacional.
	\\
	
	GitHub &
	Plataforma de hospedagem de código e colaboração científica baseada em Git.
	\\
	
	Zenodo &
	Repositório científico que permite arquivamento permanente e geração de DOI.
	\\
	
	Dryad &
	Repositório internacional para armazenamento e compartilhamento de dados científicos.
	\\
	
	Figshare &
	Plataforma de compartilhamento de dados, figuras e materiais suplementares.
	\\
	
	\bottomrule
	
\end{longtable}

\begin{conceito}
	\textbf{Observação.}
	
	A terminologia utilizada em SDM evolui continuamente à medida que novos
	algoritmos, bases de dados e abordagens metodológicas são desenvolvidos.
	Por esse motivo, recomenda-se que estudantes e pesquisadores consultem
	regularmente a literatura especializada.
\end{conceito}

%================================================
